\documentclass[11pt,a4paper]{article}
\usepackage[utf8]{inputenc}
\usepackage[T1]{fontenc}
\usepackage{lmodern}
\usepackage[french]{babel}
\frenchsetup{StandardLists=true}
\usepackage[margin=2.3cm,headheight=14pt]{geometry}
\usepackage{amsmath,amssymb}
\usepackage{icomma}
\usepackage{graphicx}
\usepackage{booktabs,array,tabularx,multirow}
\usepackage{xcolor}
\usepackage[most]{tcolorbox}
\usepackage{titlesec}
\usepackage{fancyhdr}
\usepackage{lastpage}
\usepackage{pdflscape}
\usepackage[european,straightvoltages]{circuitikz}
\usetikzlibrary{arrows.meta,positioning,fit,calc,babel}
\usepackage{enumitem}
\usepackage[hidelinks]{hyperref}

\definecolor{bleutitre}{HTML}{1F4E9A}
\definecolor{fondcalc}{HTML}{F3F6FB}
\definecolor{fonddemo}{HTML}{FBF7EE}
\definecolor{fondres}{HTML}{EEF7F1}
\definecolor{fondalerte}{HTML}{FDEDEC}

\titleformat{\section}{\Large\bfseries\color{bleutitre}}{\thesection}{0.8em}{}
\titleformat{\subsection}{\large\bfseries\color{bleutitre}}{\thesubsection}{0.8em}{}
\titleformat{\subsubsection}{\normalsize\bfseries\color{bleutitre}}{\thesubsubsection}{0.8em}{}

\newcommand{\entete}{\small Étude électrique (version 2), tour hydroponique HydroTour}
\pagestyle{fancy}
\fancyhf{}
\fancyhead[L]{\entete}
\fancyhead[R]{\small Enoaarii RAPARII}
\fancyfoot[C]{\small Page \thepage\ / \pageref{LastPage}}
\renewcommand{\headrulewidth}{0.4pt}
\fancypagestyle{plain}{\fancyhf{}\fancyhead[L]{\entete}\fancyhead[R]{\small Enoaarii RAPARII}\fancyfoot[C]{\small Page \thepage\ / \pageref{LastPage}}}

\newtcolorbox{calcul}[1]{breakable,enhanced,colback=fondcalc,colframe=bleutitre,
  fonttitle=\bfseries,title={Calcul de #1},arc=2mm,boxrule=0.6pt,before skip=10pt,after skip=10pt}
\newtcolorbox{demo}[1]{breakable,enhanced,colback=fonddemo,colframe=black!60,
  fonttitle=\bfseries,title={Démonstration : #1},arc=2mm,boxrule=0.6pt,before skip=10pt,after skip=10pt}
\newtcolorbox{resume}{breakable,enhanced,colback=fondres,colframe=black!55,
  fonttitle=\bfseries,title={Résumé},arc=2mm,boxrule=0.6pt,before skip=12pt}
\newtcolorbox{alerte}[1]{breakable,enhanced,colback=fondalerte,colframe=red!60!black,
  fonttitle=\bfseries,title={#1},arc=2mm,boxrule=0.6pt,before skip=10pt,after skip=10pt}

\newcommand{\U}[1]{\,\mathrm{#1}}
\setlength{\parskip}{4pt}
\setlength{\parindent}{0pt}
\renewcommand{\arraystretch}{1.25}
\setlength{\abovedisplayskip}{10pt}
\setlength{\belowdisplayskip}{10pt}
\setlength{\jot}{8pt}

\begin{document}

%================ PAGE DE GARDE =================
\begin{titlepage}
\thispagestyle{plain}
\centering
\vspace*{1.5cm}
{\large Université de la Polynésie française}\\[4pt]
{\large Licence 3 Sciences pour l'ingénieur (SPI)}\\[4pt]
{Projet tutoré 2026-2027}
\vspace{3.5cm}

{\color{bleutitre}\rule{\linewidth}{1pt}}\\[12pt]
{\Huge\bfseries\color{bleutitre} Étude électrique}\\[8pt]
{\Large\color{bleutitre} Version 2}\\[12pt]
{\Large Tour hydroponique verticale autonome HydroTour}\\[6pt]
{\large Schéma, bilan de consommation, dimensionnement du panneau et des batteries,\\ câbles, fusibles et protocole de mesure}\\[12pt]
{\color{bleutitre}\rule{\linewidth}{1pt}}
\vspace{3cm}

{\Large Enoaarii RAPARII}\\[6pt]
{Équipe projet de 3 étudiants}\\[6pt]
{Site de culture : UPF, Tahiti}
\vfill
{Version 2 du 1\textsuperscript{er} novembre 2026 (remplace la version 1)}
\end{titlepage}
\setcounter{page}{2}

\tableofcontents
\newpage

%================ CHANGEMENTS =================
\section*{Changements depuis la version 1}
\addcontentsline{toc}{section}{Changements depuis la version 1}

Cette version 2 reprend la structure, les données d'étiquette du panneau TE1250, l'irradiation de dimensionnement ($H = 4{,}2\U{kWh/m^2/j}$) et le rendement du convertisseur 12 V vers 5 V ($\eta = 0{,}83$) de la version 1. Les démonstrations sont réécrites et complétées. Les points suivants changent.

\begin{center}
\begin{tabularx}{\linewidth}{>{\raggedright\arraybackslash}p{3.6cm} >{\raggedright\arraybackslash}X >{\raggedright\arraybackslash}X}
\toprule
\textbf{Sujet} & \textbf{Version 1} & \textbf{Version 2} \\
\midrule
Pompes & P1 à P5, dont une pompe d'appoint P5 et un bac de secours & Bac principal unique de 152 L, appoint manuel. Pompes nommées par fonction : pompe de circulation, pompes doseuses engrais A, engrais B, pH$-$ \\
Carte de puissance & Une sortie par pompe & Carte 6318 (8 MOSFET). Circulation câblée sur 2 canaux en parallèle (0 et 1) : pertes divisées par 2 \\
Capteurs & SHT31, BH1750, ADXL345, deux eTape & DHT11, capteur de lumière analogique, LIS3DHTR (Sensor Kit), un seul eTape. ADS1115 : A0 = EC, A1 = eTape. pH sur l'entrée A1 de l'Arduino \\
Batteries & Plomb AGM supposé & LiFePO4 d'après le cahier des charges. Les deux cas sont calculés côte à côte (partie~\ref{sec:batterie}), la technologie réelle est à lire sur l'étiquette \\
Consommation & $281\U{Wh/j}$ & $226{,}1\U{Wh/j}$. Écart surtout dû à l'électronique : $91{,}8\U{Wh/j}$ estimés composant par composant, contre $149{,}8\U{Wh/j}$ en version 1 (point d'essai du convertisseur à 1~A). Pompe d'appoint retirée ; régulateur et MOSFET ajoutés. La mesure M2 tranchera \\
Production pire mois & $336\U{Wh/j}$ disponibles & $354{,}5\U{Wh/j}$ (LiFePO4) ou $317{,}2\U{Wh/j}$ (AGM), facteurs de pertes détaillés un par un \\
Fusibles & $F_{PV} = 10\U{A}$, $F_{B1} = 5\U{A}$, $F_{B2} = 15\U{A}$, $F_2 = 7{,}5\U{A}$ & $F_{PV}$, $F_{B1}$, $F_{B2}$ inchangés ; $F_2 = 4\U{A}$ (sans P5) ; nouveau $F_3 = 2\U{A}$ en entrée du convertisseur 5 V \\
Modèle du panneau & Modèle à une diode, courbes du TE1250 à 25, 47, 50 et 75~°C & Même modèle, en plus validé sur les courbes mesurées du SLP055 \\
Mesures & Absentes & Protocole complet, tableaux, CSV modèle et script de traitement (partie~\ref{sec:protocole}) \\
\bottomrule
\end{tabularx}
\end{center}

\newpage
%================ 1 =================
\section{Présentation du système électrique}

La tour est alimentée uniquement par le soleil. L'énergie suit le chemin suivant : le panneau produit, le régulateur MPPT adapte et charge les batteries, les batteries stockent pour la nuit, le convertisseur 12 V vers 5 V alimente l'électronique, et la carte MOSFET commande les pompes en 12 V.

\begin{center}
\begin{tikzpicture}[
  bloc/.style={draw=bleutitre,rounded corners=3pt,fill=fondcalc,align=center,text width=3.4cm,minimum height=2.0cm,font=\footnotesize,line width=0.8pt},
  fl/.style={-{Latex[length=2.5mm]},line width=1pt,color=black!70},
  fd/.style={-{Latex[length=2mm]},line width=0.7pt,dashed,color=bleutitre}]
\node[bloc] (pv) {\textbf{Produit l'électricité}\\ Panneau TE1250\\ 120 Wc};
\node[bloc,right=1.0cm of pv] (mppt) {\textbf{Régule et protège}\\ MPPT 20 A\\ SLC-MPPT2420D\\ fusibles};
\node[bloc,right=1.0cm of mppt] (bat) {\textbf{Stocke}\\ Batteries 12 V\\ 6 Ah + 17 Ah\\ en parallèle};
\node[bloc,below=1.2cm of mppt] (conv) {\textbf{Convertit}\\ UBEC 12 V vers 5 V\\ $\eta \approx 0{,}83$};
\node[bloc,below=1.2cm of bat] (ina) {\textbf{Mesure l'énergie}\\ INA260\\ ligne batterie};
\node[bloc,below=1.2cm of conv] (pi) {\textbf{Décide}\\ Raspberry Pi 3 B\\ + dongle Wi-Fi};
\node[bloc,left=1.0cm of pi] (ard) {\textbf{Lit et agit}\\ Arduino Uno R3\\ (USB : 5 V + données)};
\node[bloc,below=1.2cm of ard] (capt) {\textbf{Mesure l'eau et l'air}\\ pH, EC, eTape, DS18B20,\\ DHT11, lumière,\\ LIS3DHTR, ADS1115};
\node[bloc,below=1.2cm of pi] (mos) {\textbf{Commute}\\ Carte 6318\\ 8 MOSFET};
\node[bloc,right=1.0cm of mos] (pompes) {\textbf{Fait circuler et dose}\\ Pompe de circulation\\ Engrais A, engrais B, pH$-$};
\draw[fl] (pv) -- (mppt);
\draw[fl] (mppt) -- (bat);
\draw[fl] ($(mppt.south)+(-0.6,0)$) -- node[left,font=\scriptsize]{sortie LOAD} ($(conv.north)+(-0.6,0)$);
\draw[fl] (conv) -- node[left,font=\scriptsize]{5 V} (pi);
\draw[fl] (pi) -- node[above,font=\scriptsize]{USB} (ard);
\draw[fd] (ard) -- node[left,font=\scriptsize]{I2C, analogique} (capt);
\draw[fd] (ard.south east) -- node[above,font=\scriptsize,sloped]{I2C} (mos.north west);
\draw[fl] (mos) -- (pompes);
\draw[fd] (pi) -- node[above,font=\scriptsize,sloped]{I2C} (ina);
\draw[fl,color=black!45] ($(mppt.south)+(0.6,0)$) |- ($(ina.north east)+(0.5,0.4)$) |- node[right,font=\scriptsize,pos=0.25]{12 V (sortie LOAD)} (pompes.east);
\end{tikzpicture}
\end{center}

\textbf{Conventions.} Les batteries sont des batteries de 12 V nominal. Pour le LiFePO4, la tension moyenne en décharge vaut $U_B = 12{,}8\U{V}$ (4 éléments de $3{,}2\U{V}$) ; pour le plomb AGM, $U_B = 12{,}0\U{V}$. Tous les courants côté 12 V de la partie consommation sont calculés à $12{,}8\U{V}$ : à puissance égale, c'est le cas LiFePO4. Le cas AGM est traité dans la partie~\ref{sec:batterie}.

\begin{resume}
Sept fonctions : produire (panneau), réguler et protéger (MPPT, fusibles), stocker (batteries), mesurer l'énergie (INA260), convertir (UBEC), décider (Pi), lire et agir (Arduino, carte MOSFET, pompes, capteurs). Le schéma complet est donné à la page suivante.
\end{resume}

%================ 2 SCHEMA (paysage) =================
\begin{landscape}
\section{Schéma électrique complet}
\label{sec:schema}
\begin{center}
\resizebox{\linewidth}{!}{%
\begin{circuitikz}[font=\small, line width=0.7pt]
  \tikzset{boite/.style={draw=bleutitre,line width=1pt,fill=fondcalc,rounded corners=2pt,align=center}}
  % ---- rail négatif commun
  \draw[line width=1.2pt] (0,-0.4) -- (22.4,-0.4);
  \node[anchor=west,font=\footnotesize] at (22.5,-0.4) {Rail négatif commun (0 V)};
  % ---- panneau
  \draw (0,-0.4) to[pvsource, l_={\begin{tabular}{c}TE1250\\120 Wc\end{tabular}}] (0,6);
  \node[font=\footnotesize] at (-0.45,5.5) {$+$};
  \draw (0,6) to[fuse, l={$F_{PV}$ 10 A}, a={4 mm²}] (3.5,6);
  % ---- MPPT
  \draw[boite] (3.5,1) rectangle (6.5,7.5);
  \node[align=center,font=\bfseries] at (5,4.4) {Régulateur\\MPPT 20 A\\[2pt]\normalfont\footnotesize SLC-MPPT2420D};
  \node[anchor=west,font=\scriptsize] at (3.55,6) {PV$+$};
  \node[anchor=east,font=\scriptsize] at (6.45,6) {BAT$+$};
  \node[anchor=north,font=\scriptsize] at (5.6,7.5) {LOAD$+$};
  \node[anchor=south,font=\scriptsize] at (4,1) {PV$-$};
  \node[anchor=south,font=\scriptsize] at (5,1) {BAT$-$};
  \node[anchor=south,font=\scriptsize] at (6,1) {LOAD$-$};
  \draw (4,1) -- (4,-0.4) node[circ]{};
  \draw (5,1) -- (5,-0.4) node[circ]{};
  \draw (6,1) -- (6,-0.4) node[circ]{};
  % ---- INA260
  \draw[boite] (7.3,5.4) rectangle (9.5,6.6);
  \node[font=\footnotesize] at (8.4,6) {\textbf{INA260}};
  \draw (6.5,6) -- (7.3,6);
  \node[anchor=north,font=\scriptsize] at (7.7,5.4) {IN$+$};
  \node[anchor=north,font=\scriptsize] at (9.1,5.4) {IN$-$};
  \draw (8.4,5.4) -- (8.4,-0.4) node[circ]{};
  % ---- bus batteries
  \draw (9.5,6) -- (13.8,6);
  \draw (11,6) node[circ]{} to[fuse, l_={$F_{B1}$ 5 A}] (11,3.6) to[battery1, l_={\begin{tabular}{c}B1\\12 V 6 Ah\end{tabular}}, invert] (11,-0.4) node[circ]{};
  \draw (13.8,6) to[fuse, l_={$F_{B2}$ 15 A}] (13.8,3.6) to[battery1, l_={\begin{tabular}{c}B2\\12 V 17 Ah\end{tabular}}, invert] (13.8,-0.4) node[circ]{};
  \node[font=\scriptsize,align=center] at (12.4,7.0) {bus batteries, câbles 2,5 mm²\\de même longueur};
  \node[font=\scriptsize] at (8.2,6.25) {};
  % ---- ligne LOAD
  \draw (5.6,7.5) -- (5.6,8.4) -- (14.5,8.4) to[fuse, l={$F_2$ 4 A}] (17,8.4) node[circ]{};
  \node[font=\scriptsize,align=center] at (10,8.75) {sortie LOAD, 1,5 mm²};
  \node[font=\scriptsize,align=center,anchor=north west] at (17.05,8.3) {bornier 12 V};
  % ---- UBEC
  \draw (17,8.4) to[fuse, l={$F_3$ 2 A}] (19.6,8.4);
  \draw[boite] (19.6,7.8) rectangle (21.6,9.0);
  \node[font=\footnotesize,align=center] at (20.6,8.4) {\textbf{UBEC}\\12 V / 5 V};
  \draw (21.2,7.8) -- (21.2,7.2) -- (22.4,7.2) -- (22.4,-0.4) node[circ]{};
  % ---- Pi
  \draw (21.6,8.4) -- node[above,font=\scriptsize]{5 V, 0,5 mm²} (23.6,8.4);
  \draw[boite] (23.6,7.4) rectangle (27,9.4);
  \node[font=\footnotesize,align=center] at (25.3,8.4) {\textbf{Raspberry Pi 3 B}\\cerveau\\+ dongle Wi-Fi USB};
  % I2C Pi -> INA260
  \draw[dashed,bleutitre,line width=0.8pt] (25.3,9.4) -- (25.3,10.2) -- (8.4,10.2) -- (8.4,6.6);
  \node[font=\scriptsize,color=bleutitre,anchor=south] at (16,10.2) {I2C du Pi (SDA, SCL, 3,3 V) : mesure de U et I batterie};
  % ---- Arduino
  \draw[boite] (23.6,4.2) rectangle (27,6.4);
  \node[font=\footnotesize,align=center] at (25.3,5.3) {\textbf{Arduino Uno R3}\\+ base Grove\\actionneur, sécurités};
  \draw[line width=1.2pt,color=black!70] (25.3,7.4) -- node[right,font=\scriptsize,align=left]{USB\\5 V + données} (25.3,6.4);
  % ---- Capteurs
  \draw[boite] (23.6,-0.1) rectangle (27.4,3.4);
  \node[font=\scriptsize,align=left,anchor=north west] at (23.7,3.35) {\textbf{Capteurs (5 V de l'Arduino)}\\
  pH Gravity V2 : A1 Uno\\
  EC DFR0300 : ADS1115 A0\\
  \quad (alim. par D7 à la lecture)\\
  eTape 30 cm : ADS1115 A1\\
  DS18B20 : D8 \quad DHT11 : D3\\
  Lumière : A3 \quad LIS3DHTR : I2C\\
  ADS1115 0x48, RTC : I2C};
  \draw[dashed,bleutitre,line width=0.8pt] (25.3,4.2) -- (25.3,3.4);
  % ---- carte 6318
  \draw[boite] (17.8,3.2) rectangle (22.0,5.6);
  \node[font=\footnotesize,align=center] at (19.9,4.6) {\textbf{Carte Adafruit 6318}\\8 MOSFET, MCP23017 0x20};
  \draw (17,8.4) -- (17,5.0) -- (17.8,5.0);
  \node[anchor=west,font=\scriptsize] at (17.85,5.0) {V$+$};
  \draw (17.8,3.6) -- (17.4,3.6) -- (17.4,-0.4) node[circ]{};
  \node[anchor=west,font=\scriptsize] at (17.85,3.6) {GND};
  \draw[dashed,bleutitre,line width=0.8pt] (23.6,5.3) -- node[above,font=\scriptsize]{I2C} (22.0,5.3);
  % ---- pompes
  \draw (18.15,3.2) -- (18.15,2.9) -- (18.35,2.9);
  \draw (18.55,3.2) -- (18.55,2.9) -- (18.35,2.9);
  \draw (18.35,2.9) to[elmech] (18.35,1.3) -- (18.85,1.3) -- (18.85,3.2);
  \draw (19.6,3.2) to[elmech] (19.6,1.3) -- (19.95,1.3) -- (19.95,3.2);
  \draw (20.5,3.2) to[elmech] (20.5,1.3) -- (20.85,1.3) -- (20.85,3.2);
  \draw (21.4,3.2) to[elmech] (21.4,1.3) -- (21.75,1.3) -- (21.75,3.2);
  \node[font=\scriptsize,align=center] at (18.5,0.55) {circul.\\c. 0 + 1};
  \node[font=\scriptsize,align=center] at (19.75,0.55) {eng. A\\c. 2};
  \node[font=\scriptsize,align=center] at (20.65,0.55) {eng. B\\c. 3};
  \node[font=\scriptsize,align=center] at (21.55,0.55) {pH$-$\\c. 4};
\end{circuitikz}}
\end{center}

\vspace{-2mm}
{\footnotesize
\textbf{Lecture du schéma.} Trait plein : puissance. Pointillés bleus : bus de données. Chaque pompe se raccorde entre la borne V$+$ et la sortie MOSFET de son canal (commutation côté négatif, diode de roue libre intégrée à la carte). La pompe de circulation utilise les canaux 0 et 1 reliés en parallèle. L'INA260 est inséré entre la borne BAT$+$ du régulateur et le bus batteries : il mesure le courant net (charge moins décharge) et la tension batterie. Le schéma suppose que le régulateur a un négatif commun (PV$-$, BAT$-$ et LOAD$-$ reliés en interne) : à vérifier sur sa notice ; sinon, ne pas relier PV$-$ et LOAD$-$ au rail. Ordre de branchement : batteries d'abord (fusibles retirés jusqu'à la fin du câblage), panneau ensuite ; débranchement dans l'ordre inverse.
}
\end{landscape}

%================ 3 BILAN =================
\section{Bilan de consommation}
\label{sec:conso}

\subsection{Électronique alimentée en 5 V}

Le Raspberry Pi, le dongle Wi-Fi, l'Arduino (par l'USB du Pi) et tous les capteurs sont alimentés en 5 V par l'UBEC. Les valeurs sont des courants \textbf{moyens} en fonctionnement normal (hypothèses justifiées dans \texttt{HYPOTHESES.md}, à vérifier par la mesure M2).

\begin{center}
\small
\begin{tabularx}{\linewidth}{>{\raggedright\arraybackslash}X r r r}
\toprule
\textbf{Composant} & \textbf{$I$ moyen (mA)} & \textbf{$P$ à 5 V (W)} & \textbf{$E$ sur 24 h (Wh)} \\
\midrule
Raspberry Pi 3 Model B (sans écran, charge faible) & 450 & 2,250 & 54,00 \\
Dongle Wi-Fi USB à antenne externe & 120 & 0,600 & 14,40 \\
Arduino Uno R3 & 50 & 0,250 & 6,00 \\
Base Grove (LED témoin) & 3 & 0,015 & 0,36 \\
Carte pH Gravity V2.0 & 5 & 0,025 & 0,60 \\
Carte EC DFR0300 (alimentée seulement à la lecture) & 0,2 & 0,001 & 0,02 \\
eTape 30 cm et sa résistance de référence & 2 & 0,010 & 0,24 \\
ADS1115 & 0,2 & 0,001 & 0,02 \\
DS18B20 & 0,5 & 0,003 & 0,06 \\
DHT11 & 0,3 & 0,002 & 0,04 \\
Capteur de lumière (analogique) & 0,5 & 0,003 & 0,06 \\
Accéléromètre LIS3DHTR & 0,02 & 0,000 & 0,00 \\
Module RTC & 0,2 & 0,001 & 0,02 \\
INA260 (partie logique) & 1,1 & 0,006 & 0,13 \\
Carte 6318 : MCP23017 et LED témoin & 2 & 0,010 & 0,24 \\
\midrule
\textbf{Total côté 5 V} & \textbf{635} & \textbf{3,175} & \textbf{76,20} \\
\bottomrule
\end{tabularx}
\end{center}

Le Raspberry Pi et le dongle Wi-Fi représentent à eux seuls $90\U{\%}$ de la consommation de l'électronique : c'est là qu'une économie aurait un effet (désactiver le HDMI, les LED, le Bluetooth, le Wi-Fi interne si le dongle est utilisé).

\begin{demo}{énergie consommée par un composant}
La puissance instantanée reçue par un dipôle est le produit de la tension à ses bornes par le courant qui le traverse :
\[
P = U \, I
\]
L'énergie est l'intégrale de la puissance sur la durée $\Delta t$ :
\[
E = \int_{0}^{\Delta t} P(t)\,\mathrm{d}t
\]
Si $P$ est constante sur la durée de fonctionnement $\Delta t$, elle sort de l'intégrale :
\[
E = P \int_{0}^{\Delta t} \mathrm{d}t
\]
\[
E = P \, \Delta t
\]
Avec $P$ en watts et $\Delta t$ en heures, $E$ est en wattheures ($1\U{Wh} = 3600\U{J}$). Pour un courant \emph{moyen} $\bar I$ sous tension constante, $\bar P = U \bar I$ donne la même énergie que le courant réel.
\end{demo}

\begin{calcul}{l'énergie de l'électronique}
Courant total côté 5 V (somme des courants, car tous les composants sont en parallèle sur le 5 V) :
\[
I_5 = 0{,}635\U{A}
\]
Puissance côté 5 V :
\[
P_5 = U_5 \, I_5 = 5{,}0\U{V} \times 0{,}635\U{A}
\]
\[
P_5 = 3{,}175\U{W}
\]
Énergie sur une journée de 24 h :
\[
E_5 = P_5 \, \Delta t = 3{,}175\U{W} \times 24\U{h}
\]
\[
E_5 = 76{,}2\U{Wh/j}
\]
\end{calcul}

\subsection{Pertes du convertisseur 12 V vers 5 V}

\begin{demo}{puissance absorbée par un convertisseur}
Le rendement est le rapport de la puissance utile fournie à la puissance absorbée :
\[
\eta = \frac{P_{\text{sortie}}}{P_{\text{entrée}}}
\]
On isole la puissance d'entrée :
\[
P_{\text{entrée}} = \frac{P_{\text{sortie}}}{\eta}
\]
La puissance perdue (chaleur) est la différence, puis on factorise $P_{\text{sortie}}$ :
\[
P_{\text{pertes}} = P_{\text{entrée}} - P_{\text{sortie}}
\]
\[
P_{\text{pertes}} = P_{\text{sortie}} \left( \frac{1}{\eta} - 1 \right)
\]
\end{demo}

\begin{calcul}{la puissance prise sur le 12 V par l'électronique}
\[
P_{12,e} = \frac{P_5}{\eta} = \frac{3{,}175\U{W}}{0{,}83}
\]
\[
P_{12,e} = 3{,}825\U{W}
\]
\[
P_{\text{conv}} = 3{,}175\U{W} \times \left( \frac{1}{0{,}83} - 1 \right)
\]
\[
P_{\text{conv}} = 0{,}650\U{W}
\]
Courant moyen pris sur la batterie par l'électronique :
\[
I_{12,e} = \frac{P_{12,e}}{U_B} = \frac{3{,}825\U{W}}{12{,}8\U{V}}
\]
\[
I_{12,e} = 0{,}299\U{A}
\]
Énergie journalière prise sur le 12 V (dont $0{,}650\U{W} \times 24\U{h} = 15{,}6\U{Wh}$ de pertes) :
\[
E_{12,e} = 3{,}825\U{W} \times 24\U{h}
\]
\[
E_{12,e} = 91{,}8\U{Wh/j}
\]
\end{calcul}

\subsection{Pompe de circulation}

\textbf{Planning} (heure de Tahiti) : de 6 h à 18 h, 15 min de marche puis 15 min d'arrêt, soit 2 marches par heure pendant 12 h ; la nuit, 15 min à 21 h, 0 h et 3 h.

\begin{calcul}{la durée de marche journalière}
Jour : 12 heures, 2 marches de $0{,}25\U{h}$ par heure :
\[
t_{\text{jour}} = 12 \times 2 \times 0{,}25\U{h} = 6{,}00\U{h}
\]
Nuit : 3 marches de $0{,}25\U{h}$ :
\[
t_{\text{nuit}} = 3 \times 0{,}25\U{h} = 0{,}75\U{h}
\]
\[
t_c = 6{,}00\U{h} + 0{,}75\U{h} = 6{,}75\U{h/j}
\]
\end{calcul}

\begin{calcul}{l'énergie et le courant de la pompe de circulation}
\[
E_c = P_c \, t_c = 19\U{W} \times 6{,}75\U{h}
\]
\[
E_c = 128{,}25\U{Wh/j}
\]
\[
I_c = \frac{P_c}{U_B} = \frac{19\U{W}}{12{,}8\U{V}}
\]
\[
I_c = 1{,}48\U{A}
\]
\end{calcul}

\subsection{Pertes dans les MOSFET de la carte 6318}

Chaque canal de la carte utilise un transistor MOSFET dont la résistance à l'état passant vaut $R_{DS(on)} = 70\U{m\Omega}$ (valeur de fiche du composant, voir \texttt{HYPOTHESES.md}).

\begin{demo}{pertes d'un MOSFET et intérêt de deux canaux en parallèle}
Un MOSFET saturé se comporte comme une résistance $R_{DS(on)}$. La loi d'Ohm donne $U_{DS} = R_{DS(on)} I$, donc la puissance dissipée vaut :
\[
P = U_{DS} \, I = R_{DS(on)} \, I^2
\]
Avec deux canaux identiques en parallèle, les deux résistances égales se partagent le courant en deux moitiés égales :
\[
I_1 = I_2 = \frac{I}{2}
\]
La puissance totale est la somme des deux :
\[
P_{2} = R_{DS(on)} \left(\frac{I}{2}\right)^2 + R_{DS(on)} \left(\frac{I}{2}\right)^2
\]
\[
P_{2} = 2 \, R_{DS(on)} \, \frac{I^2}{4}
\]
\[
P_{2} = \frac{R_{DS(on)} \, I^2}{2}
\]
La perte totale est divisée par 2 et chaque transistor ne dissipe plus que le quart de la perte d'un canal seul : il chauffe 4 fois moins.
\end{demo}

\begin{calcul}{les pertes MOSFET de la circulation}
Un seul canal (pour comparaison) :
\[
P_1 = 0{,}070\U{\Omega} \times (1{,}48\U{A})^2 = 0{,}154\U{W}
\]
Deux canaux en parallèle :
\[
P_2 = \frac{0{,}070\U{\Omega} \times (1{,}48\U{A})^2}{2} = 0{,}077\U{W}
\]
Énergie perdue par jour :
\[
E_{\text{MOS}} = 0{,}077\U{W} \times 6{,}75\U{h} = 0{,}52\U{Wh/j}
\]
\end{calcul}

\subsection{Pompes doseuses}

Les trois pompes péristaltiques (engrais A, engrais B, pH$-$) absorbent environ $I_d = 0{,}30\U{A}$ sous 12 V et débitent environ $Q = 80\U{mL/min}$ en charge (hypothèse, à mesurer : mesure M5).

\begin{calcul}{la puissance d'une pompe doseuse}
\[
P_d = U_B \, I_d = 12{,}8\U{V} \times 0{,}30\U{A}
\]
\[
P_d = 3{,}84\U{W}
\]
\end{calcul}

\textbf{Estimation du volume d'engrais dosé par jour.} Les laitues prélèvent les sels minéraux, donc l'EC baisse. On suppose une baisse de $\Delta EC = 0{,}05\U{mS/cm}$ par jour dans le bac de $V_{\text{bac}} = 152\U{L}$, et une solution mère telle qu'ajouter $1\U{mL}$ de A et $1\U{mL}$ de B par litre d'eau augmente l'EC de $s = 0{,}56\U{mS/cm}$ (hypothèses justifiées dans \texttt{HYPOTHESES.md}, à confirmer par l'étude hydraulique et chimique).

\begin{demo}{volume de solution mère pour corriger l'EC}
En solution diluée, l'EC est proportionnelle à la concentration en sels, donc à la dose $d$ de solution mère par litre de bac :
\[
\Delta EC = s \, d
\]
La dose est le volume ajouté divisé par le volume du bac :
\[
d = \frac{V_A}{V_{\text{bac}}}
\]
On remplace et on isole $V_A$ :
\[
\Delta EC = s \, \frac{V_A}{V_{\text{bac}}}
\]
\[
V_A = \frac{\Delta EC \; V_{\text{bac}}}{s}
\]
La durée de fonctionnement d'une pompe de débit $Q$ est le volume divisé par le débit :
\[
t = \frac{V}{Q}
\]
\end{demo}

\begin{calcul}{la durée de dosage par jour}
\[
V_A = \frac{0{,}05\U{mS/cm} \times 152\U{L}}{0{,}56\U{(mS/cm)/(mL/L)}}
\]
\[
V_A = 13{,}6\U{mL/j}
\]
Le même volume d'engrais B est dosé : $V_B = 13{,}6\U{mL/j}$. Pour le pH$-$, on suppose $V_{pH} = 6\U{mL/j}$.
\[
t_A = \frac{13{,}6\U{mL}}{80\U{mL/min}} = 0{,}170\U{min} = 10{,}2\U{s}
\]
\[
t_B = 10{,}2\U{s}
\]
\[
t_{pH} = \frac{6\U{mL}}{80\U{mL/min}} = 0{,}075\U{min} = 4{,}5\U{s}
\]
Énergie estimée :
\[
E_{d,\text{est}} = 3{,}84\U{W} \times \frac{(10{,}2 + 10{,}2 + 4{,}5)\U{s}}{3600\U{s/h}}
\]
\[
E_{d,\text{est}} = 0{,}027\U{Wh/j}
\]
\textbf{Valeur retenue (majorant).} Le programme limite chaque commande à $30\U{s}$. En supposant au plus 10 commandes de $30\U{s}$ par pompe et par jour (cas d'un dérèglement) :
\[
E_{d} = 3 \times 10 \times 30\U{s} \times \frac{3{,}84\U{W}}{3600\U{s/h}}
\]
\[
E_{d} = 0{,}96\U{Wh/j}
\]
\end{calcul}

\subsection{Autoconsommation du régulateur}

\begin{calcul}{l'autoconsommation du MPPT}
Courant propre supposé $I_{\text{MPPT}} = 15\U{mA}$, jour et nuit :
\[
P_{\text{MPPT}} = 12{,}8\U{V} \times 0{,}015\U{A} = 0{,}192\U{W}
\]
\[
E_{\text{MPPT}} = 0{,}192\U{W} \times 24\U{h} = 4{,}61\U{Wh/j}
\]
\end{calcul}

\subsection{Énergie journalière totale}

\begin{center}
\begin{tabular}{l r r r}
\toprule
\textbf{Poste} & \textbf{$P$ (W)} & \textbf{Durée (h/j)} & \textbf{$E$ (Wh/j)} \\
\midrule
Électronique 5 V (vue du 12 V, pertes UBEC comprises) & 3,825 & 24 & 91,81 \\
\quad dont pertes du convertisseur & 0,650 & 24 & 15,61 \\
Pompe de circulation & 19,0 & 6,75 & 128,25 \\
Pertes MOSFET circulation (2 canaux) & 0,077 & 6,75 & 0,52 \\
Pompes doseuses (majorant) & 3,84 & 0,25 & 0,96 \\
Autoconsommation du MPPT & 0,192 & 24 & 4,61 \\
\midrule
\textbf{Total} & & & \textbf{226,1} \\
\bottomrule
\end{tabular}
\end{center}

\begin{calcul}{l'énergie journalière}
\[
E_j = 91{,}81 + 128{,}25 + 0{,}52 + 0{,}96 + 4{,}61\U{Wh/j}
\]
\[
E_j = 226{,}1\U{Wh/j}
\]
La pompe de circulation représente $128{,}25 / 226{,}1 = 57\U{\%}$ du besoin, l'électronique $41\U{\%}$.
\end{calcul}

\begin{resume}
Besoin journalier : $E_j = 226{,}1\U{Wh/j}$ (version 1 : $281\U{Wh/j}$). Circulation $128{,}25\U{Wh}$ (57 \%), électronique $91{,}8\U{Wh}$ dont $15{,}6\U{Wh}$ de pertes du convertisseur, doseuses moins de $1\U{Wh}$. Courant de repos sur le 12 V : environ $0{,}31\U{A}$ ; courant pompe de circulation : $1{,}48\U{A}$. Les deux canaux MOSFET en parallèle divisent les pertes par 2 ($0{,}077\U{W}$ au lieu de $0{,}154\U{W}$).
\end{resume}

%================ 4 PANNEAU =================
\newpage
\section{Dimensionnement du panneau}
\label{sec:panneau}

\subsection{Données d'étiquette du TE1250}

\begin{center}
\begin{tabular}{l c c}
\toprule
\textbf{Grandeur (conditions STC : 1000 W/m², 25 °C)} & \textbf{Symbole} & \textbf{Valeur} \\
\midrule
Puissance crête & $P_c$ & $120\U{Wc}$ \\
Puissance minimale garantie & $P_{\min}$ & $117{,}5\U{W}$ \\
Tension au point de puissance maximale & $V_{mp}$ & $17{,}6\U{V}$ \\
Courant au point de puissance maximale & $I_{mp}$ & $6{,}8\U{A}$ \\
Courant de court-circuit & $I_{sc}$ & $7{,}3\U{A}$ \\
Tension de circuit ouvert & $V_{oc}$ & $22{,}0\U{V}$ \\
Calibre maximal du fusible série & & $10\U{A}$ \\
\bottomrule
\end{tabular}
\end{center}

Les coefficients de température du TE1250 ne figurent pas sur l'étiquette : on prend ceux du SLP055 (même technologie 36 cellules) : $\beta_{V} = -80\U{mV/°C}$, $\alpha_{I} = +0{,}065\U{\%/°C}$, $\gamma_P = -0{,}5\U{\%/°C}$, NOCT $= 47\U{°C}$.

\subsection{Modèle à une diode du panneau}

\begin{demo}{équation du modèle à une diode}
\textbf{Une cellule.} Une cellule éclairée se représente par :
\begin{itemize}[nosep]
\item une source de courant $I_{ph}$ (photocourant), proportionnelle à l'éclairement ;
\item une diode (la jonction PN) traversée par le courant $I_D$ ;
\item une résistance parallèle $R_{sh}$ (fuites), traversée par $I_{sh}$ ;
\item une résistance série $R_s$ (contacts, grilles, rubans), traversée par le courant de sortie $I$.
\end{itemize}
\textbf{Loi des nœuds} au point commun de la source, de la diode et de $R_{sh}$ :
\[
I_{ph} = I_D + I_{sh} + I
\]
\[
I = I_{ph} - I_D - I_{sh}
\]
\textbf{Tension interne.} La diode et $R_{sh}$ sont soumises à la tension de sortie augmentée de la chute dans $R_s$ (loi des mailles) :
\[
V_j = V + R_s I
\]
\textbf{Diode.} Le courant d'une jonction PN suit la loi exponentielle :
\[
I_D = I_0 \left[ \exp\!\left( \frac{V_j}{n V_T} \right) - 1 \right]
\]
avec $I_0$ le courant de saturation, $n$ le facteur d'idéalité (entre 1 et 2) et $V_T$ la tension thermique :
\[
V_T = \frac{k \, T}{q}
\]
($k = 1{,}381 \times 10^{-23}\U{J/K}$, $q = 1{,}602 \times 10^{-19}\U{C}$, $T$ en kelvins ; à $25\U{°C}$, $V_T = 25{,}7\U{mV}$).

\textbf{Résistance parallèle} (loi d'Ohm) :
\[
I_{sh} = \frac{V_j}{R_{sh}}
\]
\textbf{Panneau de $N_s$ cellules en série.} Le même courant traverse toutes les cellules et leurs tensions s'additionnent : on remplace $n V_T$ par $a = n N_s V_T$ et on utilise les résistances totales $R_s$ et $R_{sh}$ du panneau. En remplaçant $I_D$ et $I_{sh}$ dans la loi des nœuds :
\[
I = I_{ph} - I_0 \left[ \exp\!\left( \frac{V + R_s I}{a} \right) - 1 \right] - \frac{V + R_s I}{R_{sh}}
\]
$I$ figure des deux côtés : l'équation est implicite et se résout numériquement (recherche du zéro pour chaque tension $V$).
\end{demo}

\begin{demo}{identification des paramètres à partir de l'étiquette}
\textbf{Court-circuit} ($V = 0$, $I = I_{sc}$). La tension interne $R_s I_{sc}$ est très petite, la diode ne conduit presque pas :
\[
I_{sc} \approx I_{ph} - \frac{R_s I_{sc}}{R_{sh}}
\]
\[
I_{ph} = I_{sc} \left( 1 + \frac{R_s}{R_{sh}} \right)
\]
\textbf{Circuit ouvert} ($I = 0$, $V = V_{oc}$) :
\[
0 = I_{ph} - I_0 \left[ \exp\!\left( \frac{V_{oc}}{a} \right) - 1 \right] - \frac{V_{oc}}{R_{sh}}
\]
\[
I_0 = \frac{I_{ph} - V_{oc}/R_{sh}}{\exp\left( V_{oc}/a \right) - 1}
\]
\textbf{Point de puissance maximale.} Il reste deux inconnues, $n$ et $R_s$ ($R_{sh}$ est fixée à $200\U{\Omega}$, valeur courante, peu influente). Deux conditions les déterminent : la courbe passe par $(V_{mp}, I_{mp})$, et la puissance y est maximale, donc sa dérivée est nulle :
\[
I(V_{mp}) = I_{mp}
\]
\[
\left. \frac{\mathrm{d}(V I)}{\mathrm{d}V} \right|_{V_{mp}} = 0
\]
Ces deux équations sont résolues numériquement (script \texttt{modele\_panneau.py}).

\textbf{Effet de la température.} D'après les coefficients, pour une température de cellule $T_c$ :
\[
I_{sc}(T_c) = I_{sc,\text{STC}} \left[ 1 + \alpha_I (T_c - 25\U{°C}) \right]
\]
\[
V_{oc}(T_c) = V_{oc,\text{STC}} + \beta_V (T_c - 25\U{°C})
\]
$I_{ph}$ et $I_0$ sont recalculés avec ces valeurs, et $a$ avec la nouvelle $V_T$.
\end{demo}

\begin{calcul}{les paramètres du TE1250}
Résultat de la résolution numérique :
\[
n = 1{,}063 \qquad R_s = 0{,}234\U{\Omega} \qquad R_{sh} = 200\U{\Omega}
\]
\[
a = n \, N_s \, V_T = 1{,}063 \times 36 \times 0{,}0257\U{V} = 0{,}983\U{V}
\]
\[
I_{ph} = 7{,}3\U{A} \times \left( 1 + \frac{0{,}234\U{\Omega}}{200\U{\Omega}} \right) = 7{,}309\U{A}
\]
\[
I_0 = \frac{7{,}309\U{A} - 22{,}0\U{V}/200\U{\Omega}}{\exp(22{,}0\U{V}/0{,}983\U{V}) - 1} = 1{,}37 \times 10^{-9}\U{A}
\]
\end{calcul}

\textbf{Validation du modèle sur le SLP055.} La même méthode appliquée au SLP055 ($n = 1{,}74$, $R_s = 0{,}143\U{\Omega}$) donne les courbes de la figure~\ref{fig:valid}, comparées aux points de puissance maximale lus sur ses courbes fournies.

\begin{figure}[h]
\centering
\includegraphics[width=0.85\linewidth]{figures/validation_slp055.pdf}
\caption{Modèle à une diode du SLP055 et points de puissance maximale lus sur ses courbes.}
\label{fig:valid}
\end{figure}

\begin{center}
\begin{tabular}{c c c c}
\toprule
\textbf{Température} & \textbf{$P_{\max}$ lue (W)} & \textbf{$P_{\max}$ modèle (W)} & \textbf{Écart} \\
\midrule
25 °C & 55,0 & 55,0 & 0,0 \% \\
47 °C & 50,1 & 49,3 & $-$1,5 \% \\
50 °C & 49,5 & 48,6 & $-$1,9 \% \\
75 °C & 43,8 & 42,1 & $-$3,9 \% \\
\bottomrule
\end{tabular}
\end{center}

Le modèle est légèrement pessimiste à chaud (moins de 4 \% d'écart) : il convient pour dimensionner.

\begin{figure}[p]
\centering
\includegraphics[width=0.92\linewidth]{figures/te1250_courbes.pdf}
\caption{Caractéristiques I(U) et P(U) du TE1250 calculées par le modèle à une diode, sous 1000 W/m².}
\label{fig:te1250}
\end{figure}

\begin{center}
\begin{tabular}{c c c c c}
\toprule
\textbf{$T_c$} & \textbf{$V_{mp}$ (V)} & \textbf{$I_{mp}$ (A)} & \textbf{$P_{\max}$ (W)} & \textbf{$V_{oc}$ (V)} \\
\midrule
25 °C & 17,60 & 6,80 & 119,7 & 22,00 \\
47 °C & 15,82 & 6,83 & 107,9 & 20,24 \\
50 °C & 15,58 & 6,83 & 106,3 & 20,00 \\
54 °C & 15,26 & 6,83 & 104,2 & 19,68 \\
75 °C & 13,60 & 6,82 & 92,8 & 18,00 \\
\bottomrule
\end{tabular}
\end{center}

Même à $75\U{°C}$, $V_{mp} = 13{,}6\U{V}$ reste au-dessus de la tension de charge d'une batterie 12 V vide : le régulateur MPPT peut toujours charger. C'est la principale raison d'utiliser un régulateur MPPT plutôt qu'une simple liaison directe.

\subsection{Température des cellules}

\begin{demo}{température de cellule à partir de la NOCT}
La cellule reçoit une puissance proportionnelle à l'éclairement $G$ et la cède à l'air par convection et rayonnement, de façon à peu près proportionnelle à l'écart de température. En régime établi, puissance reçue $=$ puissance cédée :
\[
\alpha \, G = h \, (T_c - T_a)
\]
\[
T_c - T_a = \frac{\alpha}{h} \, G = K \, G
\]
La NOCT est la température atteinte sous $G = 800\U{W/m^2}$ avec $T_a = 20\U{°C}$ :
\[
\text{NOCT} - 20\U{°C} = K \times 800\U{W/m^2}
\]
\[
K = \frac{\text{NOCT} - 20\U{°C}}{800\U{W/m^2}}
\]
\[
T_c = T_a + \frac{\text{NOCT} - 20\U{°C}}{800\U{W/m^2}} \; G
\]
\end{demo}

\begin{calcul}{la température de cellule en milieu de journée (pire mois)}
Hypothèses : $T_a = 27\U{°C}$ et $G = 800\U{W/m^2}$ en milieu de journée.
\[
T_c = 27\U{°C} + \frac{47\U{°C} - 20\U{°C}}{800\U{W/m^2}} \times 800\U{W/m^2}
\]
\[
T_c = 54\U{°C}
\]
\end{calcul}

\begin{calcul}{le facteur de température}
Avec le coefficient de puissance $\gamma_P = -0{,}5\U{\%/°C}$ :
\[
k_T = 1 + \gamma_P \, (T_c - 25\U{°C})
\]
\[
k_T = 1 - 0{,}005\U{/°C} \times (54\U{°C} - 25\U{°C})
\]
\[
k_T = 0{,}855
\]
\textbf{Vérification sur les courbes du SLP055.} Interpolation linéaire entre $50\U{°C}$ ($49{,}5\U{W}$) et $75\U{°C}$ ($43{,}8\U{W}$) :
\[
P(54\U{°C}) = 49{,}5\U{W} + (43{,}8 - 49{,}5)\U{W} \times \frac{54 - 50}{75 - 50}
\]
\[
P(54\U{°C}) = 48{,}6\U{W}
\]
\[
\frac{48{,}6\U{W}}{55{,}0\U{W}} = 0{,}883
\]
Les courbes donnent une perte un peu plus faible ($0{,}883$) que le coefficient ($0{,}855$). On garde $k_T = 0{,}855$, qui est plus prudent. Le modèle du TE1250 donne $104{,}2 / 119{,}7 = 0{,}870$, entre les deux.
\end{calcul}

\subsection{Énergie produite le pire mois}

\begin{demo}{énergie produite à partir de l'irradiation journalière}
La puissance d'un panneau est proportionnelle à l'éclairement $G$ ; à $G_{\text{STC}} = 1000\U{W/m^2}$ elle vaut $P_c$ :
\[
P(t) = P_c \, \frac{G(t)}{G_{\text{STC}}}
\]
On intègre sur la journée :
\[
E_{\text{th}} = \int_{\text{jour}} P(t) \, \mathrm{d}t = \frac{P_c}{G_{\text{STC}}} \int_{\text{jour}} G(t)\,\mathrm{d}t
\]
L'intégrale de l'éclairement sur la journée est l'irradiation journalière $H$ :
\[
E_{\text{th}} = P_c \, \frac{H}{G_{\text{STC}}}
\]
Le rapport $H / G_{\text{STC}}$ est un nombre d'heures, appelé heures équivalentes de plein soleil. Les pertes réelles sont prises en compte par un produit de facteurs (chaque perte s'applique à ce qui reste après la précédente) :
\[
E_{\text{prod}} = E_{\text{th}} \; k_T \; k_{\text{tol}} \; k_s \; k_{\text{MPPT}} \; k_{\text{câbles}} \; k_{\text{bat}}
\]
\end{demo}

\textbf{Facteurs retenus.}
\begin{itemize}[nosep]
\item $k_T = 0{,}855$ : température des cellules (calcul précédent) ;
\item $k_{\text{tol}} = 117{,}5 / 120 = 0{,}979$ : tolérance de fabrication (puissance minimale garantie de l'étiquette) ;
\item $k_s = 0{,}95$ : salissures, sel marin, poussière, orientation imparfaite ;
\item $k_{\text{MPPT}} = 0{,}95$ : rendement du régulateur ;
\item $k_{\text{câbles}} = 0{,}98$ : chutes de tension dans les câbles (partie~\ref{sec:cables}) ;
\item $k_{\text{bat}} = 0{,}95$ (LiFePO4) ou $0{,}85$ (AGM) : rendement énergétique charge puis décharge. Il est appliqué à toute l'énergie, alors qu'une partie est consommée directement le jour : c'est prudent.
\end{itemize}

\begin{calcul}{l'énergie produite par le TE1250 le pire mois}
Irradiation du pire mois retenue : $H = 4{,}2\U{kWh/m^2/j}$.
\[
E_{\text{th}} = 120\U{W} \times \frac{4200\U{Wh/m^2/j}}{1000\U{W/m^2}}
\]
\[
E_{\text{th}} = 120\U{W} \times 4{,}2\U{h/j} = 504\U{Wh/j}
\]
Rapport de performance (cas LiFePO4) :
\[
PR = 0{,}855 \times 0{,}979 \times 0{,}95 \times 0{,}95 \times 0{,}98 \times 0{,}95
\]
\[
PR = 0{,}703
\]
\[
E_{\text{prod}} = 504\U{Wh/j} \times 0{,}703
\]
\[
E_{\text{prod}} = 354{,}5\U{Wh/j}
\]
Cas AGM ($k_{\text{bat}} = 0{,}85$) : $PR = 0{,}629$ et $E_{\text{prod}} = 317{,}2\U{Wh/j}$.
\end{calcul}

\begin{calcul}{la marge du panneau}
\[
m = \frac{E_{\text{prod}}}{E_j} - 1
\]
\[
m_{\text{LFP}} = \frac{354{,}5\U{Wh/j}}{226{,}1\U{Wh/j}} - 1 = 0{,}57 = +57\U{\%}
\]
\[
m_{\text{AGM}} = \frac{317{,}2\U{Wh/j}}{226{,}1\U{Wh/j}} - 1 = 0{,}40 = +40\U{\%}
\]
\textbf{Sensibilité.} Avec une irradiation plus faible de $H = 3{,}9\U{kWh/m^2/j}$, la production LiFePO4 tombe à $329{,}2\U{Wh/j}$ : la marge reste de $+46\U{\%}$.
\end{calcul}

\textbf{Comparaison avec l'ancien panneau SLP055.} Avec les mêmes facteurs (hors tolérance), il produirait $55\U{W} \times 4{,}2\U{h} \times 0{,}718 = 166\U{Wh/j}$, soit moins que le besoin de $226{,}1\U{Wh/j}$ : il ne peut servir que d'appoint.

\begin{resume}
Le modèle à une diode, validé à moins de 4 \% sur les courbes du SLP055, donne pour le TE1250 : $P_{\max} = 104{,}2\U{W}$ à $15{,}3\U{V}$ pour une cellule à $54\U{°C}$. Le pire mois ($H = 4{,}2\U{kWh/m^2/j}$), le panneau fournit $354{,}5\U{Wh/j}$ en LiFePO4 ($317{,}2\U{Wh/j}$ en AGM) pour un besoin de $226{,}1\U{Wh/j}$ : marge de $+57\U{\%}$ ($+40\U{\%}$ en AGM). Le panneau de 120 Wc est suffisant ; le SLP055 seul ne l'est pas.
\end{resume}

%================ 5 BATTERIE =================
\newpage
\section{Dimensionnement des batteries}
\label{sec:batterie}

\subsection{LiFePO4 ou plomb AGM : les deux cas}

Le cahier des charges indique des batteries LiFePO4, la version 1 les supposait au plomb AGM. Les deux technologies se comportent très différemment ; le tableau donne les valeurs retenues pour chacune.

\begin{center}
\begin{tabular}{l c c}
\toprule
\textbf{Grandeur} & \textbf{LiFePO4} & \textbf{Plomb AGM} \\
\midrule
Tension nominale de calcul $U_B$ & 12,8 V & 12,0 V \\
Profondeur de décharge admise (DoD) & 80 \% & 50 \% \\
Rendement énergétique charge et décharge & 0,95 & 0,85 \\
Courant de charge admis (usage courant) & $0{,}5\,C$ & $0{,}25\,C$ \\
Tension de charge (absorption) & 14,2 à 14,6 V & 14,4 à 14,7 V \\
Tension de maintien (floating) & aucune ou 13,6 V & 13,5 à 13,8 V \\
Tension au repos, batterie pleine & environ 13,4 V & environ 12,8 V \\
Tension au repos, batterie vide & environ 12,0 V & environ 11,9 V (à 50 \%) \\
Protection interne (BMS) & oui, en général & non \\
\midrule
Énergie nominale (23 Ah) & 294,4 Wh & 276,0 Wh \\
Énergie utilisable (23 Ah) & 235,5 Wh & 138,0 Wh \\
\bottomrule
\end{tabular}
\end{center}

\begin{alerte}{À vérifier par l'étudiant avant tout raccordement}
\begin{itemize}[nosep]
\item Lire sur l'\textbf{étiquette de chaque batterie} sa technologie (LiFePO4, \og Lithium Iron Phosphate \fg, ou \og SLA, AGM, Lead Acid \fg).
\item \textbf{Mettre en parallèle deux technologies différentes est interdit} : tensions de charge, courbes de décharge et protections différentes ; l'une serait surchargée ou détruite. Si les deux batteries ne sont pas de la même technologie, n'en utiliser qu'une, ou en remplacer une.
\item Régler le régulateur MPPT sur le type de batterie réel (profil lithium ou plomb).
\end{itemize}
\end{alerte}

\begin{demo}{énergie utilisable d'une batterie}
Une capacité $C$ (en ampères-heures) délivrée sous une tension moyenne $U_B$ correspond à l'énergie (même raisonnement que $E = P \Delta t$ avec $P = U_B I$) :
\[
E_{\text{nom}} = U_B \, C
\]
On ne s'autorise à prélever que la fraction DoD pour préserver la durée de vie :
\[
E_u = U_B \, C \; \text{DoD}
\]
Deux batteries de même tension en parallèle additionnent leurs capacités (les courants s'additionnent à tension égale) :
\[
C = C_1 + C_2 = 6\U{Ah} + 17\U{Ah} = 23\U{Ah}
\]
\end{demo}

\begin{calcul}{l'énergie utilisable}
\[
E_{u,\text{LFP}} = 12{,}8\U{V} \times 23\U{Ah} \times 0{,}80 = 235{,}5\U{Wh}
\]
\[
E_{u,\text{AGM}} = 12{,}0\U{V} \times 23\U{Ah} \times 0{,}50 = 138{,}0\U{Wh}
\]
\end{calcul}

\subsection{Autonomie de nuit}

La nuit compte 14 h sans production utile (de 17 h à 7 h, hypothèse prudente pour juin). Pendant cette période, la circulation tourne 7 fois 15 min (17 h, 17 h 30, 21 h, 0 h, 3 h, 6 h, 6 h 30).

\begin{calcul}{l'énergie consommée la nuit}
\[
t_{c,\text{nuit}} = 7 \times 0{,}25\U{h} = 1{,}75\U{h}
\]
\[
E_n = (3{,}825 + 0{,}192)\U{W} \times 14\U{h} + (19 + 0{,}077)\U{W} \times 1{,}75\U{h}
\]
\[
E_n = 56{,}24\U{Wh} + 33{,}38\U{Wh} = 89{,}6\U{Wh}
\]
Part de l'énergie utilisable :
\[
\frac{E_n}{E_{u,\text{LFP}}} = \frac{89{,}6\U{Wh}}{235{,}5\U{Wh}} = 38\U{\%}
\]
\[
\frac{E_n}{E_{u,\text{AGM}}} = \frac{89{,}6\U{Wh}}{138{,}0\U{Wh}} = 65\U{\%}
\]
Profondeur de décharge réelle au matin : $89{,}6 / 294{,}4 = 30\U{\%}$ (LiFePO4) ou $89{,}6 / 276{,}0 = 32\U{\%}$ (AGM).
\end{calcul}

\subsection{Jours sans soleil}

\begin{calcul}{l'autonomie sans aucune production}
\[
N_0 = \frac{E_u}{E_j}
\]
\[
N_{0,\text{LFP}} = \frac{235{,}5\U{Wh}}{226{,}1\U{Wh/j}} = 1{,}04\U{j}
\]
\[
N_{0,\text{AGM}} = \frac{138{,}0\U{Wh}}{226{,}1\U{Wh/j}} = 0{,}61\U{j}
\]
\end{calcul}

\begin{calcul}{l'autonomie par jours très couverts}
Un jour très couvert apporte encore environ $H_c = 1{,}2\U{kWh/m^2/j}$ (hypothèse). Production :
\[
E_{c,\text{LFP}} = 120\U{W} \times 1{,}2\U{h} \times 0{,}703 = 101{,}3\U{Wh/j}
\]
Déficit journalier :
\[
D_{\text{LFP}} = 226{,}1\U{Wh/j} - 101{,}3\U{Wh/j} = 124{,}9\U{Wh/j}
\]
Nombre de jours couverts consécutifs supportés en partant de batteries pleines :
\[
N_{c,\text{LFP}} = \frac{235{,}5\U{Wh}}{124{,}9\U{Wh/j}} = 1{,}9\U{j}
\]
Cas AGM : $E_c = 90{,}6\U{Wh/j}$, $D = 135{,}5\U{Wh/j}$, $N_c = 138{,}0 / 135{,}5 = 1{,}0\U{j}$.
\end{calcul}

\textbf{Interprétation.} Les batteries couvrent largement la nuit, mais pas une longue période couverte. Le programme du Raspberry Pi doit donc prévoir un \textbf{mode sobriété} déclenché sur la tension batterie mesurée par l'INA260 (par exemple sous $12{,}9\U{V}$ en LiFePO4 au repos) : circulation réduite à 10 min par heure le jour, nuit inchangée. La consommation descend alors vers $154\U{Wh/j}$. La batterie de 23 Ah pour deux jours sans soleil complet demanderait :
\[
C_{2j} = \frac{2 \times 226{,}1\U{Wh}}{12{,}8\U{V} \times 0{,}80} = 44\U{Ah} \quad \text{(LiFePO4)}
\]
\[
C_{2j} = \frac{2 \times 226{,}1\U{Wh}}{12{,}0\U{V} \times 0{,}50} = 75\U{Ah} \quad \text{(AGM)}
\]

\subsection{Mise en parallèle de deux batteries de capacités différentes}

\begin{demo}{partage du courant entre deux batteries en parallèle}
Chaque batterie se modélise par une force électromotrice $E_k$ en série avec une résistance interne $R_k$. En parallèle, elles ont la même tension $U$ aux bornes :
\[
U = E_1 - R_1 I_1 = E_2 - R_2 I_2
\]
Si les deux batteries sont au même état de charge, $E_1 = E_2$, donc :
\[
R_1 I_1 = R_2 I_2
\]
Avec $I = I_1 + I_2$ le courant total :
\[
I_1 = I \, \frac{R_2}{R_1 + R_2}
\qquad
I_2 = I \, \frac{R_1}{R_1 + R_2}
\]
Pour une même technologie, la résistance interne est à peu près inversement proportionnelle à la capacité, $R_k = r / C_k$. On remplace :
\[
I_1 = I \, \frac{r/C_2}{r/C_1 + r/C_2}
\]
On multiplie haut et bas par $C_1 C_2 / r$ :
\[
I_1 = I \, \frac{C_1}{C_1 + C_2}
\]
Chaque batterie prend une part du courant proportionnelle à sa capacité : elles se chargent et se déchargent au même régime ($I_k / C_k$ identique).

\textbf{Courant de circulation.} Si l'on relie deux batteries dont les tensions diffèrent de $\Delta E = E_1 - E_2$, sans charge extérieure, la maille donne :
\[
I_{\text{circ}} = \frac{E_1 - E_2}{R_1 + R_2}
\]
\end{demo}

\begin{calcul}{le partage du courant de charge et le courant de circulation}
Courant de charge maximal sortant du régulateur :
\[
I_{\text{ch}} = \frac{P_c \; k_{\text{MPPT}}}{U_B} = \frac{120\U{W} \times 0{,}95}{12{,}8\U{V}} = 8{,}9\U{A}
\]
\[
I_{B1} = 8{,}9\U{A} \times \frac{6\U{Ah}}{23\U{Ah}} = 2{,}32\U{A} \quad (0{,}39\,C)
\]
\[
I_{B2} = 8{,}9\U{A} \times \frac{17\U{Ah}}{23\U{Ah}} = 6{,}58\U{A} \quad (0{,}39\,C)
\]
En LiFePO4, $0{,}39\,C < 0{,}5\,C$ : acceptable. En AGM, $0{,}39\,C > 0{,}25\,C$ : charge un peu trop forte en plein soleil (vieillissement accéléré).

Courant de circulation si l'on relie les batteries avec $0{,}3\U{V}$ d'écart ($R_{B1} = 60\U{m\Omega}$ supposé, donc $R_{B2} = 60 \times 6/17 = 21\U{m\Omega}$) :
\[
I_{\text{circ}} = \frac{0{,}3\U{V}}{0{,}060\U{\Omega} + 0{,}021\U{\Omega}} = 3{,}7\U{A}
\]
Un faible écart de tension provoque déjà un courant important : d'où les précautions suivantes.
\end{calcul}

\textbf{Précautions de mise en parallèle.}
\begin{enumerate}[nosep]
\item Même technologie et même tension nominale (obligatoire).
\item Charger chaque batterie séparément jusqu'à pleine charge, puis vérifier au multimètre un écart de tension inférieur à $0{,}05\U{V}$ avant de les relier.
\item Un fusible par batterie ($F_{B1}$, $F_{B2}$), au plus près de la borne $+$ : un court-circuit interne d'une batterie ne doit pas être alimenté par l'autre.
\item Câbles de même section et de même longueur entre chaque batterie et le bus.
\item Raccordement \og en diagonale \fg : le $+$ du bus pris sur une batterie et le $-$ sur l'autre, pour équilibrer les chutes de tension.
\item En LiFePO4, chaque batterie a son BMS : si l'un coupe, toute la charge passe dans l'autre (voir le choix des fusibles, partie~\ref{sec:cables}).
\item Remplacer les deux batteries ensemble en fin de vie.
\end{enumerate}

\begin{resume}
Capacité totale $23\U{Ah}$. Énergie utilisable : $235{,}5\U{Wh}$ en LiFePO4 (DoD 80 \%), $138{,}0\U{Wh}$ en AGM (DoD 50 \%). La nuit consomme $89{,}6\U{Wh}$, soit 38 \% (LiFePO4) ou 65 \% (AGM) de l'énergie utilisable. Autonomie sans soleil : $1{,}0\U{j}$ (LiFePO4) ou $0{,}6\U{j}$ (AGM) ; par jours très couverts : $1{,}9\U{j}$ ou $1{,}0\U{j}$. Un mode sobriété piloté par la tension batterie est nécessaire. La technologie réelle des deux batteries doit être lue sur leurs étiquettes : jamais deux technologies différentes en parallèle.
\end{resume}

%================ 6 CABLES FUSIBLES =================
\newpage
\section{Sections de câbles et fusibles}
\label{sec:cables}

\subsection{Chute de tension et section minimale}

\begin{demo}{section minimale d'un câble}
La résistance d'un conducteur de longueur $\ell$ et de section $S$ vaut :
\[
R = \rho \, \frac{\ell}{S}
\]
Le courant fait l'aller et le retour : pour une liaison de longueur $L$, $\ell = 2L$ :
\[
R = \frac{2 \rho L}{S}
\]
La chute de tension (loi d'Ohm) :
\[
\Delta U = R \, I = \frac{2 \rho L I}{S}
\]
On impose $\Delta U \le \varepsilon \, U$ (chute relative maximale $\varepsilon$) :
\[
\frac{2 \rho L I}{S} \le \varepsilon U
\]
\[
S \ge \frac{2 \rho L I}{\varepsilon U}
\]
La résistivité du cuivre augmente avec la température :
\[
\rho(T) = \rho_{20} \left[ 1 + \alpha_{Cu} (T - 20\U{°C}) \right]
\]
\end{demo}

\begin{calcul}{la résistivité du cuivre à 60 °C}
Câbles au soleil ou dans le boîtier chaud : $T = 60\U{°C}$.
\[
\rho = 0{,}0172\U{\Omega\,mm^2/m} \times \left[ 1 + 0{,}00393\U{/°C} \times (60 - 20)\U{°C} \right]
\]
\[
\rho = 0{,}0199\U{\Omega\,mm^2/m}
\]
\end{calcul}

\begin{calcul}{la section du câble panneau vers régulateur}
$L = 5\U{m}$, $I = 6{,}83\U{A}$ et $U = 15{,}26\U{V}$ (point de puissance maximale à $54\U{°C}$), $\varepsilon = 3\U{\%}$ :
\[
S \ge \frac{2 \times 0{,}0199\U{\Omega\,mm^2/m} \times 5\U{m} \times 6{,}83\U{A}}{0{,}03 \times 15{,}26\U{V}}
\]
\[
S \ge \frac{1{,}359\U{\Omega\,mm^2\,A}}{0{,}458\U{V}} = 2{,}97\U{mm^2}
\]
On choisit $4\U{mm^2}$ (câble solaire). Chute réelle :
\[
\Delta U = \frac{2 \times 0{,}0199 \times 5 \times 6{,}83}{4}\U{V} = 0{,}34\U{V} \quad (2{,}2\U{\%})
\]
\end{calcul}

Les autres liaisons sont calculées de la même façon :

\begin{center}
\small
\begin{tabular}{l c c c c c c c}
\toprule
\textbf{Liaison} & \textbf{$L$ (m)} & \textbf{$I$ (A)} & \textbf{$U$ (V)} & \textbf{$\varepsilon$} & \textbf{$S_{\min}$ (mm²)} & \textbf{$S$ (mm²)} & \textbf{$\Delta U$ réel} \\
\midrule
Panneau vers MPPT & 5,0 & 6,83 & 15,3 & 3 \% & 2,97 & 4 & 2,2 \% \\
MPPT vers INA260 vers bus & 0,5 & 8,91 & 12,8 & 1 \% & 1,38 & 4 & 0,35 \% \\
Bus vers B1 (6 Ah) & 0,3 & 2,32 & 12,8 & 0,5 \% & 0,43 & 2,5 & 0,09 \% \\
Bus vers B2 (17 Ah) & 0,3 & 6,58 & 12,8 & 0,5 \% & 1,23 & 2,5 & 0,25 \% \\
Sortie LOAD vers bornier & 0,5 & 2,70 & 12,8 & 1 \% & 0,42 & 1,5 & 0,28 \% \\
Bornier vers circulation & 3,0 & 1,48 & 12,8 & 3 \% & 0,46 & 1,5 & 0,92 \% \\
Bornier vers doseuse & 1,0 & 0,30 & 12,8 & 3 \% & 0,03 & 0,5 & 0,19 \% \\
UBEC vers Pi (5 V) & 0,3 & 1,00 & 5,0 & 2 \% & 0,12 & 0,5 & 0,48 \% \\
\bottomrule
\end{tabular}
\end{center}

La chute totale du panneau aux batteries vaut environ $2{,}2 + 0{,}35 = 2{,}6\U{\%}$ au courant maximal, et nettement moins en moyenne sur la journée : le facteur $k_{\text{câbles}} = 0{,}98$ est cohérent. Sur le 5 V du Pi, une chute supérieure à $0{,}25\U{V}$ provoque des alertes de sous-tension : le câble court de $0{,}5\U{mm^2}$ est indispensable (éviter un câble USB fin et long).

\subsection{Fusibles}

\begin{demo}{règle de choix d'un fusible}
Le fusible doit laisser passer le courant normal maximal $I_b$ sans fondre ni vieillir, d'où une marge de $25\U{\%}$ :
\[
I_n \ge 1{,}25 \, I_b
\]
et il doit fondre avant que le câble protégé chauffe dangereusement, donc rester sous le courant admissible du câble $I_z$ :
\[
1{,}25 \, I_b \le I_n \le I_z
\]
On prend le calibre normalisé immédiatement supérieur à $1{,}25\,I_b$ (valeurs courantes : 1, 2, 3, 4, 5, 7,5, 10, 15, 20 A).
\end{demo}

\begin{calcul}{les calibres des fusibles}
\textbf{$F_{PV}$ (panneau).} Le courant maximal est le courant de court-circuit, augmenté à chaud :
\[
1{,}25 \times 7{,}3\U{A} = 9{,}1\U{A} \;\Rightarrow\; I_n = 10\U{A}
\]
C'est aussi le calibre maximal autorisé par l'étiquette du panneau.

\textbf{$F_{B1}$ (batterie 6 Ah).} Courant normal en charge $2{,}32\U{A}$ :
\[
1{,}25 \times 2{,}32\U{A} = 2{,}9\U{A} \;\Rightarrow\; I_n = 5\U{A} \text{ (conservé)}
\]
Si le BMS de B2 coupe, les $8{,}9\U{A}$ de charge passeraient dans B1 : le fusible de 5 A fond et protège la petite batterie.

\textbf{$F_{B2}$ (batterie 17 Ah).} Normalement $6{,}58\U{A}$, mais jusqu'à $8{,}9\U{A}$ si B1 est déconnectée :
\[
1{,}25 \times 8{,}9\U{A} = 11{,}1\U{A} \;\Rightarrow\; I_n = 15\U{A} \text{ (conservé)}
\]

\textbf{$F_2$ (sortie LOAD, toutes les charges).} Courant maximal : circulation, trois doseuses ensemble (cas extrême) et électronique :
\[
I_b = 1{,}48\U{A} + 3 \times 0{,}30\U{A} + 0{,}30\U{A} = 2{,}68\U{A}
\]
\[
1{,}25 \times 2{,}68\U{A} = 3{,}4\U{A} \;\Rightarrow\; I_n = 4\U{A}
\]
(La version 1 prévoyait $7{,}5\U{A}$ avec la pompe d'appoint ; sans elle, $4\U{A}$ protège mieux.)

\textbf{$F_3$ (entrée UBEC).} Courant d'entrée maximal si l'UBEC fournit ses $3\U{A}$ nominaux :
\[
I_{\text{entrée,max}} = \frac{5\U{V} \times 3\U{A}}{0{,}83 \times 12{,}8\U{V}} = 1{,}41\U{A}
\]
\[
1{,}25 \times 1{,}41\U{A} = 1{,}8\U{A} \;\Rightarrow\; I_n = 2\U{A}
\]
\end{calcul}

\begin{center}
\begin{tabular}{l l c c c}
\toprule
\textbf{Fusible} & \textbf{Emplacement} & \textbf{$I_b$ (A)} & \textbf{Calibre v1} & \textbf{Calibre v2} \\
\midrule
$F_{PV}$ & panneau, borne $+$ & 7,3 & 10 A & 10 A \\
$F_{B1}$ & batterie 6 Ah, borne $+$ & 2,32 & 5 A & 5 A \\
$F_{B2}$ & batterie 17 Ah, borne $+$ & 6,58 (8,9) & 15 A & 15 A \\
$F_2$ & sortie LOAD du MPPT & 2,68 & 7,5 A & 4 A \\
$F_3$ & entrée de l'UBEC & 0,30 (1,41) & absent & 2 A \\
\bottomrule
\end{tabular}
\end{center}

\textbf{Autres vérifications.} L'INA260 supporte $15\U{A}$ en continu : avec $F_{B2} = 15\U{A}$ en amont et $8{,}9\U{A}$ maximum en fonctionnement, il est protégé. La tension maximale du panneau (matin frais, $T_c = 20\U{°C}$) vaut $22{,}0 - 0{,}080 \times (20 - 25) = 22{,}4\U{V}$ : à vérifier inférieure à la tension PV maximale de la notice du régulateur. Chaque canal de la carte 6318 supporte $3{,}6\U{A}$ en pointe : la circulation ($0{,}74\U{A}$ par canal) et les doseuses ($0{,}30\U{A}$) sont loin de cette limite.

\begin{resume}
Câbles : $4\U{mm^2}$ panneau vers régulateur (chute $2{,}2\U{\%}$) et régulateur vers bus, $2{,}5\U{mm^2}$ vers chaque batterie, $1{,}5\U{mm^2}$ pour la sortie LOAD et la pompe de circulation, $0{,}5\U{mm^2}$ pour les doseuses et le 5 V du Pi. Fusibles : $F_{PV} = 10\U{A}$, $F_{B1} = 5\U{A}$, $F_{B2} = 15\U{A}$ (inchangés), $F_2 = 4\U{A}$ (au lieu de 5 A), nouveau $F_3 = 2\U{A}$ en entrée de l'UBEC.
\end{resume}

%================ 7 PROTOCOLE =================
\newpage
\section{Protocole de mesure}
\label{sec:protocole}

Ce protocole permet de remplacer les hypothèses de cette étude par des valeurs mesurées. Les fichiers sont dans le dossier \texttt{protocole/} :
\begin{itemize}[nosep]
\item \texttt{modele\_mesures\_ponctuelles.csv} : modèle vide pour les mesures M1, M2, M4, M5 ;
\item \texttt{modele\_mesures\_energie.csv} : modèle vide pour l'enregistrement d'une journée (M3) ;
\item \texttt{traitement\_mesures.py} : script de traitement ; lancé sans argument, il fonctionne sur des données d'\textbf{EXEMPLE} fictives (\texttt{exemple\_*.csv}).
\end{itemize}

\textbf{Matériel.} Multimètre avec calibre 10 A continu ; INA260 déjà câblé (lecture par le Raspberry Pi) ; testeur USB (courant 5 V) si disponible ; éprouvette graduée de 100 mL et chronomètre ; thermomètre infrarouge (facultatif) ; fiches de relevé ci-dessous.

\textbf{Sécurité.} Ne jamais brancher le multimètre en ampèremètre en parallèle sur la batterie (court-circuit). Ouvrir le circuit, brancher le multimètre en série, puis refermer. Retirer le fusible concerné pendant les manipulations de câbles.

\subsection{M1 : courant de chaque pompe}
\begin{enumerate}[nosep]
\item Retirer $F_2$, débrancher le fil $+$ de la pompe sur la carte 6318, insérer le multimètre (calibre 10 A) en série.
\item Remettre $F_2$, commander la pompe depuis le Pi (\texttt{D 0 60000} pour la circulation, \texttt{D 1 10000} pour l'engrais A, etc.).
\item Relever le courant après 5 s de marche (régime établi) et la tension aux bornes de la pompe.
\item Recommencer pour chaque pompe ; pour les doseuses, le tuyau doit être en charge (plongé dans le liquide).
\end{enumerate}

\begin{center}
\begin{tabular}{|p{4.2cm}|p{2.2cm}|p{2.2cm}|p{2.2cm}|p{2.8cm}|}
\hline
\textbf{Pompe} & \textbf{$U$ mesurée (V)} & \textbf{$I$ mesuré (A)} & \textbf{$P = UI$ (W)} & \textbf{Théorie} \\
\hline
Circulation (canaux 0 + 1) & & & & 1,48 A, 19 W \\[6pt]\hline
Doseuse engrais A & & & & 0,30 A \\[6pt]\hline
Doseuse engrais B & & & & 0,30 A \\[6pt]\hline
Doseuse pH$-$ & & & & 0,30 A \\[6pt]\hline
\end{tabular}
\end{center}

\subsection{M2 : consommation au repos et rendement de l'UBEC}
\begin{enumerate}[nosep]
\item La nuit (pas de charge), toutes pompes arrêtées, le Pi en fonctionnement normal.
\item Relever le courant batterie donné par l'INA260 pendant 10 min (une valeur par minute) et en faire la moyenne.
\item Contrôler une valeur au multimètre en série sur la sortie LOAD.
\item Si un testeur USB est disponible, mesurer le courant 5 V du Pi ; calculer le rendement $\eta = (U_5 I_5)/(U_{12} I_{12})$.
\end{enumerate}

\begin{center}
\begin{tabular}{|p{5cm}|p{3cm}|p{3cm}|p{3cm}|}
\hline
\textbf{Grandeur} & \textbf{Mesure} & \textbf{Unité} & \textbf{Théorie} \\
\hline
Courant de repos 12 V (INA260, moyenne) & & A & 0,314 \\[6pt]\hline
Courant de repos 12 V (multimètre) & & A & 0,314 \\[6pt]\hline
Tension de sortie UBEC & & V & 5,0 \\[6pt]\hline
Courant 5 V (testeur USB) & & A & 0,635 \\[6pt]\hline
Rendement UBEC calculé & & sans unité & 0,83 \\[6pt]\hline
\end{tabular}
\end{center}

\subsection{M3 : production et bilan sur une journée}
\begin{enumerate}[nosep]
\item Choisir une journée représentative ; noter la météo (ensoleillé, nuageux, pluie).
\item Enregistrer toutes les 5 min, de 0 h à 24 h, la tension et le courant batterie de l'INA260 (le cerveau peut l'écrire directement dans le CSV ; courant positif en charge).
\item De 6 h à 18 h, relever toutes les 30 min la tension et le courant du panneau (afficheur du régulateur ou multimètre), et mettre 0 la nuit. Toutes les lignes doivent être remplies.
\item Noter l'état de la pompe de circulation (1 marche, 0 arrêt).
\item Lancer \texttt{python3 traitement\_mesures.py journee.csv ponctuel.csv}.
\end{enumerate}

\begin{center}
\begin{tabular}{|p{1.4cm}|p{2.0cm}|p{2.0cm}|p{2.0cm}|p{2.0cm}|p{1.4cm}|p{2.0cm}|}
\hline
\textbf{Heure} & \textbf{$U_{\text{bat}}$ (V)} & \textbf{$I_{\text{bat}}$ (A)} & \textbf{$I_{PV}$ (A)} & \textbf{$U_{PV}$ (V)} & \textbf{Pompe} & \textbf{Remarque} \\
\hline
06:00 & & & & & & \\\hline
08:00 & & & & & & \\\hline
10:00 & & & & & & \\\hline
12:00 & & & & & & \\\hline
14:00 & & & & & & \\\hline
16:00 & & & & & & \\\hline
18:00 & & & & & & \\\hline
\end{tabular}
\end{center}

\subsection{M4 : tension batterie et identification}
\begin{enumerate}[nosep]
\item Lire et noter l'étiquette de chaque batterie (technologie, capacité, tension de charge, courant maximal).
\item Avant de les relier, mesurer la tension à vide de chaque batterie au multimètre.
\item Mesurer la tension du bus au repos avant l'aube (fin de nuit), puis à midi en charge, et pendant une marche de la pompe.
\item Comparer chaque valeur à celle de l'INA260 (écart attendu inférieur à $0{,}05\U{V}$).
\end{enumerate}

\begin{center}
\begin{tabular}{|p{5.2cm}|p{2.6cm}|p{2.6cm}|p{3.6cm}|}
\hline
\textbf{Situation} & \textbf{Multimètre (V)} & \textbf{INA260 (V)} & \textbf{Repère LiFePO4} \\
\hline
Technologie B1 / B2 (étiquette) & \multicolumn{2}{l|}{} & identiques ? \\[6pt]\hline
B1 seule, à vide & & & 13,0 à 13,4 \\[6pt]\hline
B2 seule, à vide & & & 13,0 à 13,4 \\[6pt]\hline
Bus, fin de nuit (repos) & & & $\ge 13{,}0$ \\[6pt]\hline
Bus, midi en charge & & & 13,8 à 14,6 \\[6pt]\hline
Bus, pompe en marche & & & \\[6pt]\hline
\end{tabular}
\end{center}

\subsection{M5 : débit des pompes doseuses}
\begin{enumerate}[nosep]
\item Plonger le tuyau d'aspiration dans de l'eau, placer la sortie dans l'éprouvette.
\item Commander la pompe 30 s (\texttt{D 1 30000}) et lire le volume.
\item Débit $Q = V / t$ ; refaire 3 fois et faire la moyenne.
\end{enumerate}

\begin{center}
\begin{tabular}{|p{3.6cm}|p{2.4cm}|p{2.4cm}|p{2.4cm}|p{3cm}|}
\hline
\textbf{Pompe} & \textbf{Essai 1 (mL)} & \textbf{Essai 2 (mL)} & \textbf{Essai 3 (mL)} & \textbf{$Q$ moyen (mL/min)} \\
\hline
Engrais A & & & & \\[6pt]\hline
Engrais B & & & & \\[6pt]\hline
pH$-$ & & & & \\[6pt]\hline
\end{tabular}
\end{center}

\subsection{Traitement : énergie par la méthode des trapèzes}

\begin{demo}{méthode des trapèzes}
L'énergie est l'intégrale de la puissance $p(t) = u(t)\,i(t)$. On ne connaît $p$ qu'aux instants de relevé $t_0, t_1, \dots, t_N$. Entre deux relevés, on remplace $p$ par le segment de droite qui joint $(t_k, p_k)$ et $(t_{k+1}, p_{k+1})$ :
\[
p(t) \approx p_k + \frac{p_{k+1} - p_k}{t_{k+1} - t_k} \, (t - t_k)
\]
On intègre cette droite sur $[t_k, t_{k+1}]$, en posant $h = t_{k+1} - t_k$ et $\tau = t - t_k$ :
\[
\int_{t_k}^{t_{k+1}} p(t)\,\mathrm{d}t \approx \int_{0}^{h} \left( p_k + \frac{p_{k+1} - p_k}{h} \, \tau \right) \mathrm{d}\tau
\]
\[
= p_k \, h + \frac{p_{k+1} - p_k}{h} \cdot \frac{h^2}{2}
\]
\[
= h \, \frac{p_k + p_{k+1}}{2}
\]
C'est l'aire d'un trapèze de bases $p_k$ et $p_{k+1}$ et de hauteur $h$. On additionne tous les intervalles :
\[
E \approx \sum_{k=0}^{N-1} (t_{k+1} - t_k) \, \frac{p_k + p_{k+1}}{2}
\]
Avec $t$ en heures et $p$ en watts, $E$ est en wattheures.
\end{demo}

Le script calcule ainsi l'énergie produite par le panneau $E_{PV}$ et l'énergie nette entrée dans la batterie $E_{\text{net}}$, puis la consommation journalière par le bilan :
\[
E_{\text{cons}} = k_{\text{MPPT}} \, E_{PV} - E_{\text{net}}
\]
(ce que le régulateur fournit, moins ce qui reste stocké dans la batterie). Il compare $E_{PV}$ à $354{,}5\U{Wh/j}$, $E_{\text{cons}}$ à $226{,}1\U{Wh/j}$ et chaque mesure ponctuelle à sa valeur théorique.

\textbf{Résultat sur les données d'EXEMPLE (fictives, pour preuve de fonctionnement).}

\begin{center}
\small
\begin{tabular}{l r r r}
\toprule
\textbf{Grandeur (EXEMPLE)} & \textbf{Calcul} & \textbf{Théorie} & \textbf{Écart} \\
\midrule
Énergie produite par le panneau & 350,5 Wh/j & 354,5 Wh/j & $-$1,1 \% \\
Énergie nette entrée dans la batterie & 107,8 Wh/j & & \\
Consommation déduite du bilan & 225,2 Wh/j & 226,1 Wh/j & $-$0,4 \% \\
Courant de la pompe de circulation & 1,52 A & 1,48 A & $+$2,4 \% \\
Courant de repos 12 V & 0,33 A & 0,314 A & $+$5,1 \% \\
\bottomrule
\end{tabular}
\end{center}

\begin{figure}[h]
\centering
\includegraphics[width=0.9\linewidth]{protocole/exemple_bilan_journee.png}
\caption{Sortie du script sur les données d'EXEMPLE : puissances du panneau et de la batterie, tension batterie.}
\end{figure}

\begin{figure}[h]
\centering
\includegraphics[width=0.85\linewidth]{protocole/exemple_ecarts.png}
\caption{Sortie du script sur les données d'EXEMPLE : écart entre mesures ponctuelles et théorie.}
\end{figure}

\textbf{Exploitation.} Si un écart dépasse $10\U{\%}$, remplacer l'hypothèse correspondante par la valeur mesurée dans \texttt{calculs.py}, relancer le calcul et mettre à jour le bilan de la partie~\ref{sec:conso}.

\begin{resume}
Cinq séries de mesures : courant de chaque pompe (M1), consommation au repos et rendement de l'UBEC (M2), bilan d'une journée enregistré par l'INA260 toutes les 5 min (M3), tensions et identification des batteries (M4), débit des doseuses (M5). Le script intègre les puissances par la méthode des trapèzes, déduit la consommation réelle et la compare à la théorie ; il a été vérifié sur des données d'exemple.
\end{resume}

%================ 8 SYNTHESE =================
\section{Synthèse}

\begin{center}
\begin{tabular}{l c c}
\toprule
\textbf{Résultat} & \textbf{Version 1} & \textbf{Version 2} \\
\midrule
Consommation journalière & 281 Wh/j & 226,1 Wh/j \\
Production théorique (120 Wc, $H = 4{,}2$) & 504 Wh/j & 504 Wh/j \\
Production disponible pire mois & 336 Wh/j & 354,5 (LFP) / 317,2 (AGM) Wh/j \\
Marge du panneau & $+20$ \% & $+57$ \% (LFP) / $+40$ \% (AGM) \\
Énergie utilisable batteries (23 Ah) & & 235,5 (LFP) / 138,0 (AGM) Wh \\
Part de la nuit dans l'énergie utilisable & & 38 \% (LFP) / 65 \% (AGM) \\
Autonomie sans soleil & & 1,0 j (LFP) / 0,6 j (AGM) \\
Fusibles & 10 / 5 / 15 / 7,5 A & 10 / 5 / 15 / 4 A + 2 A (UBEC) \\
\bottomrule
\end{tabular}
\end{center}

\begin{resume}
Le panneau TE1250 de 120 Wc couvre le besoin du pire mois avec une marge confortable. Les batteries de 23 Ah couvrent la nuit sans difficulté mais seulement un à deux jours très couverts en LiFePO4 (un seul en AGM) : un mode sobriété basé sur la tension batterie est à programmer. Points à vérifier en priorité par l'étudiant : technologie réelle des deux batteries (étiquette), réglage du régulateur, négatif commun du régulateur, puis les mesures M1 à M5.
\end{resume}

\end{document}
